% !TeX root = ../papers/nc-p3.tex
\section{反対八次曲面の核像の次数と二十平面を除いた曲線}
\label{sec:octic-kernel-image-degree}
\begin{theorem}\label{thm:octic-kernel-image-degree}
四辺有限の標準有限Frobenius窓で $D=V(\det N)$ は整、$X$ は入力長さ20の点窓とする。
階数2以下の軌跡が有限なら、一般核の閉包 $Z$ は曲面で、$D\dashrightarrow Z$ の次数を $n$ として
\[
               (\deg Z,n)=(3,4),(4,3),(6,2),(12,1)
\]
のいずれかである。出力一般平面による切断には算術種数25の整曲線から次数 $n$ の有限射がある。
十五次元成分の稠密開集合では四位置とも $(12,1)$ で、一般平面切断の構造層からこの有限双有理像の
押出し構造層への余核の長さは30である。
\end{theorem}
\begin{proof}[三十二の交点と垂直成分の重複度]
核の入射 $Y=\{(x,v):N(x)v=0\}$ は $\mathbb P^3\times\mathbb P^3$ で四本の $(2,1)$ 式で定まる。
$D$ の階数3部分では核のグラフであり、階数2以下が有限なら全体が期待次元2の完全交叉となる。
定理\ref{thm:double-point-opposite-reconstruction}で三次小行列式の層は $\mathcal I_X^2$ に等しい。
従って特殊台は $X$ の台に限られ、各台点 $p$ の上の平面 $F_p$ は出力点 $y(p)$ の直交平面である。
主成分は $D$ と双有理で一般点で被約である。

$F_p$ の重複度は $\ell(\mathcal O_{X,p})$ である。実際 $p$ の周囲で単元一成分を消し、残る対称三行を $S$ とする。
$I=\mathcal I_{X,p}$ は三元正則列イデアルであり、$S$ の成分は $I$ に入る。
三次小行列式の等号により $I_2(S)=I^2$ となる。$I/\mathfrak m_p I$ の三基底を変数 $u$ として
一次項を $S_1(u)$ とする。六つの二次小行列式は $I^2/\mathfrak m_p I^2=\operatorname{Sym}^2(I/\mathfrak m_p I)$ を張るので、
$u\ne0$ で $S_1(u)$ の階数は2以上である。$\det S_1$ は恒等零ではない。
さもなくば平面上の全ての形式が相異なる二直線の積となり、二因子を順序付けるエタール二重被覆は分裂する。
二射の次数の和が1なので固定因子と一次で動く因子となり、動く因子は射影平面全体を走って平方を生じる。
従って階数2となる方向は高々曲線で、各方向の核は一点、一般 $w\in\mathbb P^2$ に対して
$u\mapsto S_1(u)w$ は同型である。局所入射で一つの出力座標を消去した後の三式は $S(x)w=0$ である。
一般 $w$ で中山の補題によりこれらは $I$ 自体を生成する。これで垂直重複度が判明する。

出力の一般直線を引き戻すと全交点数は $(2h+k)^4 k^2=32h^3k^3$ である。
垂直平面は合計20を与えるから主成分の寄与は12となり、$Z$ は曲面で $n\deg Z=12$ である。
十個の六次余因子は同じ定理の極小生成元なので独立であり、非退化な $Z$ を含む二次曲面は存在しない。
これで表の候補を得る。出力全体での四つの二次式の結果式から垂直因子を重複度付きで除いたものは、
$Z$ の方程式の $n$ 乗となる。これは結果式の因子が完全交叉入射の射影サイクルであり、
一般直線における各局所交点長を重複度に取ることによる。
\end{proof}

\begin{proof}[一般平面と残余の算術種数]
出力を一般平面に制限した曲線 $Y_H$ は $\mathbb P^3\times\mathbb P^2$ の四本の $(2,1)$ の完全交叉である。
その双対化層は $\mathcal O(4,1)$、次数は $(4h+k)(2h+k)^4=128h^3k^2$ で、$\chi(\mathcal O_{Y_H})=-64$。
点窓 $X$ の上で $y\cdot v=0$ を課した $V\subset Y_H$ は $X$ 上の射影直線束である。
平面はどの $F_p$ とも異なるためである。$V$ は局所完全交叉で、$Y_H$ 内の残余 $B$ について
\[
 0\longrightarrow\omega_V\otimes\omega_{Y_H}^{-1}\longrightarrow\mathcal O_{Y_H}
       \longrightarrow\mathcal O_B\longrightarrow0
\]
となる。左端は $V$ のイデアルの零化イデアルであり、Gorenstein双対性による。
各局所Artin成分上の直線束で左端は $\mathcal O(-3)$ だからEuler数は $-2\ell(X)=-40$、
従って $p_a(B)=25$ である。前段の局所計算により垂直直線は $V$ に全重複度で入り、残余にその成分はない。
主成分の一般平面切断は像が二次元なのでBertiniにより既約で一般点で被約である。
完全交叉内の局所完全交叉の残余はCohen--Macaulayなので $B$ は整である。
主成分の出力射影の正次元繊維は高々有限個の出力点上に限られ、一般の出力平面はこれらを避ける。
よって $B\to Z\cap H$ は有限で次数 $n$ である。$n=1$ のとき平面十二次式の算術種数55との
差が二つの構造層の余核長 $55-25=30$ を与える。
\end{proof}

\begin{proof}[十五次元成分での次数の決定]
系\ref{cor:normal-octics-general-finite-chart}の有理標本を用いる。
\path{computations/verify_octic_kernel_image_degree.py} は良い素数2657で四位置の二次式行列を復元し、
出力直線 $v=(2,17,11,5)+s(31,7,19,29)$ 上で $N(x)v$ の四二次式のMacaulay結果式を取る。
次数5の単項式で分子を作り、二種類以上の平方で割れる単項式の部分行列式を分母に取る。
分母が非零の33値で次数32の式を補間し、更に二値で照合する。対応する二十の出力一次因子を除くと、
各位置で次数12、その導関数との最大公約因子は1となる。従って標数零の同じ標本の残余結果式も平方因子を持たない。
$n>1$ ならどの直線でも制限が零か重根を持つため、四位置とも $n=1$ である。
被約な点窓をエタール局所的に標識すれば垂直因子の積も代数的に変化し、この非消滅条件は開である。
既知の正規性の開条件と合わせると当該成分の稠密開集合で同じ結論となる。
\end{proof}

\begin{corollary}\label{cor:excess-kernel-degree-eleven}
階数1軌跡 $R$ のみが有限で階数2の曲線がある場合、$Z$ が曲面なら $n\deg Z\le11$ である。
特に九余因子で二次曲面を支配すれば $n\le5$ となる。
\end{corollary}
\begin{proof}
今も全入射は完全交叉で、垂直平面の重複度の和は $\ell(R)\ge20$ であるか、それ以上である。
実際各一般平面点で $S(x)w$ のイデアルは $S$ の全成分のイデアルに入る。
従って $n\deg Z\le12$ で、等号なら $R=X$ であり、階数2曲線の上の核直線族は全て出力次数零である。
その直線が動けば出力像は曲面となるので、各曲線で核二空間は固定となる。
入力曲線上の $N$ は固定三空間 $\operatorname{Sym}^2 k^2$ に入る。全二次式を使う写像なので
この入力曲線は直線である（非共線な曲線の二次評価は少なくとも五次元）。
その直線で $N$ は対称二行二次式行列となり、階数1切断は非零の四次行列式で切られ長さ4である。
$R=X$ と直線切断長3以下に反する。九余因子なら有限階数2軌跡は既に排除されているので結論を得る。
\end{proof}
